% !TeX root = main.tex
% ============================================================================
% preamble.tex -- everything before \begin{document}: packages, class options, the ABNT/UFV
% metadata fields, hyperref setup and the build switches' consumers.
%
% WHY THIS FILE EXISTS (author decision 2026-07-29). This and content.tex were one file, 0_main.tex,
% which was 44 percent preamble and 56 percent document. The name described neither half, and two
% things depended on the boundary between them:
%   * src_utils/mkformat.py precompiles the preamble into a format dump. It found the boundary by
%     searching for \begin{document} and hashing the bytes before it, so its staleness key contained
%     a BYTE OFFSET that moved on every preamble edit.
%   * A reader had to know that the top of a 580-line file was preamble and the rest was not.
% Splitting makes both names true and removes the offset from the staleness key: the format now
% depends on preamble.tex's own checksum.
%
% WHAT BELONGS HERE: anything that must be read before \begin{document}. If you add a package or a
% metadata field, it goes here and the format dump is invalidated automatically -- `make fast3` will
% rebuild it. Anything that produces output belongs in content.tex.
%
% DO NOT put \begin{document} in this file. mkformat.py dumps everything here and stops; a
% \begin{document} would be swallowed into the format and every later build would fail with
% "Extra }, or forgotten \endgroup" or a silently empty document.
% ============================================================================

% !TeX root = main.tex
%%==============================================================
%% Master's dissertation -- Vitor H. O. Silva, UFV / PPGCC (Campus Florestal)
%% Skeleton derived from a prior UFV/PPGCC abnTeX2 dissertation tree (TEMPLATE.md §0
%% decision), machinery kept, content stripped.
%% KEPT: abntex2 class + abntex2-UFV.sty + abntex2-num.bst, memoir options,
%%       hyperref setup, babel plumbing (brazil loaded, english main), booktabs/
%%       subcaption/amsmath float stack, frenchspacing.
%% STRIPPED: the source tree's chapters, \modelname macro, color palette, listings/
%%       algorithm2e/lipsum/svg/pdfpages/dsfont/lscape (unused here; re-add per need),
%%       and its Abstract text and title metadata.
%% CHANGED: lmodern -> newtxtext/newtxmath (Times; manual §8);
%%       abntex2cite [alf] -> [num] (settled decision #5, single numeric global list);
%%       two build modes via \ifdefensebuild (TEMPLATE.md §2 item 4);
%%       chapterpreface environment (time-capsule device, NORTH_STAR §3);
%%       \OnehalfSpacing enforced (manual §7: 1.5 line spacing).
%% Compile with pdflatex via main.tex (defense), main_ppgc.tex (defense + approval sheet),
%% or `make final` (AcademicoPG deposit body). Never compile this file directly.
%%
%% ONE THING TO KNOW BEFORE YOU EDIT THE PREAMBLE (everything from here down to
%% \begin{document}). A precompiled pdflatex FORMAT DUMP of this preamble exists, built by
%% src_utils/mkformat.py into build/fmt/mainpre.fmt and used by the `make fast` / `make fast3`
%% targets. A format dump is a frozen snapshot of the preamble: if it is loaded after the
%% preamble has changed, pdflatex uses the OLD preamble, the build succeeds, and the PDF is
%% silently wrong -- no warning, no error, nothing in the log that names the cause. That is the
%% worst failure shape this repository has (science/AGENT_HANDOFF.md §2.3b and §2.4b: a clean
%% report over a wrong artifact).
%%
%% YOU DO NOT HAVE TO REMEMBER TO REBUILD IT. mkformat.py keys the dump to a checksum of this
%% preamble plus main.tex and the files they load, rebuilds when the key moves, and
%% src_utils/fastbuild.sh REFUSES to run against a stale key rather than producing a build from
%% it (`fastbuild: REFUSING to build -- stale: ...`). The plain `make defense`/`final`/`ppgc`
%% targets never load the dump at all, so they are always correct and are what Overleaf uses.
%% If a `fast` build ever disagrees with a plain build, the plain build is the document: report
%% it, do not work around it. `make format-status` says whether the dump is fresh and why.
%%==============================================================

% [round6, E-5] hyperfootnotes=false must reach hyperref AT LOAD TIME. abntex2 loads hyperref
% itself, and \hypersetup after the fact does not undo the footnote-anchor mechanism -- setting
% it there left all ten "name{Hfootnote.N} has been referenced but does not exist" warnings in
% place, measured. \PassOptionsToPackage before \documentclass is the only hook that reaches it.
% See the fuller note at the \hypersetup block below for what the defect was.
\PassOptionsToPackage{hyperfootnotes=false}{hyperref}

\documentclass[
    12pt,                    % manual §7: size 12
    openright,
    oneside,
    a4paper,
    sumario=tradicional,
    brazil,                    % PT available for Resumo and the folha de rosto boilerplate
    english,                % last = main document language
]{abntex2}

% -- Machinery (from the source tree, kept) --
\usepackage{abntex2-UFV}        % UFV front matter + 3/2 cm margins (checked vs manual §7)
% Gate D-1 fix: abntex2-UFV runs \checkandfixthelayout under single spacing, leaving a
% text block ~11pt taller than the 3cm/2cm margins dictate (bottom margin measured
% 1.5-1.6 cm vs the 2 cm spec). Re-derive the block from the margins under the 1.5 line
% spacing the manual mandates, then fix WITHOUT line-rounding.
%
% THE FOUR LINES BELOW ARE ONE UNIT AND THEIR ORDER IS LOAD-BEARING. \OnehalfSpacing FIRST,
% because \checkandfixthelayout derives the text block from the CURRENT \baselineskip; the
% two \set...marginsandblock lines declare the margins UFV manual §7 specifies; and the
% [fixed] option on the last line is not decoration -- it is the reason the bottom margin
% comes out at the specified 2 cm.
%
% WHY [fixed] IS LOAD-BEARING, MEASURED 2026-07-29 on this preamble's own package stack
% (a minimal document with these five lines, compiled from src/ with only the option
% changed; both variants tex_errors=0):
%   \checkandfixthelayout[fixed]   \textheight 702.78308pt   \lowermargin 56.9055pt = 2.0000 cm
%   \checkandfixthelayout          \textheight 713.78398pt   \lowermargin 45.9046pt = 1.6134 cm
% The default option is `classic`, and it does NOT keep the height the margins dictate. It
% REPLACES it with a whole number of lines plus \topskip: memoir.cls:1039 \m@mclassicht
% computes floor(\textheight / \baselineskip) * \baselineskip + \topskip. Here that is
% floor(702.78308 / 17.99446) = 39 lines, so 39 * 17.99446 + 12pt = 713.78394pt -- 11.00pt
% TALLER than the margins asked for. A taller block with the top margin held fixed can only
% grow downward, so the whole 11.00pt is taken out of the bottom margin: 2.0000 cm becomes
% 1.6134 cm, against a specification of 2 cm. [fixed] is the one option value that performs
% no line-rounding at all (memoir.cls:1096-1137: `classic`, `lines` and `nearest` each call
% a rounding macro; `fixed` calls none), which is why it is the option here.
%
% HOW THE BREAKAGE PRESENTS, which is the reason this is worth a comment: it does not.
% pdflatex raises no warning, no overfull box, no error; every gate in src_utils/check.sh
% stays green; \OnehalfSpacing is still in force so the text looks entirely normal. The
% document simply prints with a bottom margin ~4 mm short of the specification on every
% page, and the reader who finds it is whoever checks the deposit against the UFV manual.
% If you touch these four lines, re-derive \lowermargin the way the numbers above were
% obtained rather than trusting that the build stayed clean.
\OnehalfSpacing
\setlrmarginsandblock{3cm}{2cm}{*}
\setulmarginsandblock{3cm}{2cm}{*}
\checkandfixthelayout[fixed]
\usepackage[T1]{fontenc}
\usepackage[utf8]{inputenc}
\usepackage{newtxtext,newtxmath} % Times New Roman equivalent (manual §8); replaces lmodern
\usepackage{indentfirst}
\usepackage{graphicx}
\usepackage{adjustbox}
\usepackage{booktabs}
\usepackage{multirow}
\usepackage{multicol}
\usepackage{amsmath} % amssymb/amsfonts dropped: newtxmath provides the AMS symbols (clash on \Bbbk)
\usepackage{mathtools}
\usepackage{subcaption}
\usepackage{caption}
\usepackage{xurl}
\DisemulatePackage{setspace}
\usepackage{setspace}
\usepackage{bm}
\usepackage{longtable}

% -- Citations: single global NUMERIC scheme (decision #5) --
\usepackage[num,abnt-emphasize=bf]{abntex2cite}
% [round5, advisor request 2026-07-27] Citation brackets: square, not round.
% His words: "sugiro mudar a notacao das citacoes tambem. Esta com (NUMERO) e na listagem esta com o
% numero sem []. Acho melhor usar o [NUMERO] e quando necessario, colocar o nome dos autores com et al."
% abntex2cite exposes \citebrackets{open}{close}; the num style calls \setcitebrackets which defaults
% to (). Overriding it AFTER the package is loaded is the supported hook (abntex2cite.sty:212-213).
% Precedent: the Viegas dissertation, whose structure this document follows, renders 193 [N] citations
% against 15 (N); Germano, canesche, passe and lapsusvgi all use (N). So both are attested in the
% program and this is a style choice, which the advisor has now made.
\citebrackets{[}{]}
% The reference LIST label is built separately, by \@biblabel via \citenumstyle
% (abntex2cite.sty:834-848). \citenumstyle is a zero-argument FONT SWITCH, not a formatter, so
% redefining it to take the number renders [0]; \@biblabel is the right hook. Redefined AFTER
% \begin{document} would be too late, and the package sets it inside its own \AtBeginDocument, so
% this hooks in at the same point and wins by ordering.
\makeatletter
\AtBeginDocument{\renewcommand{\@biblabel}[1]{%
    \ifthenelse{\equal{#1}{}}{}{{\citenumstyle[#1]\hspace{\biblabelsep}}}}}
\makeatother

% -- Babel caption plumbing (from the source tree, kept: brazil loaded via class option,
%    caption names anglicized, then english selected as the running language) --
\addto\captionsbrazil{%
    \def\bibname{References}%
    \renewcommand{\figurename}{Figure}
    \renewcommand{\tablename}{Table}
    \renewcommand{\listfigurename}{List of Figures}
    \renewcommand{\listtablename}{List of Tables}
    \renewcommand{\contentsname}{Contents}
    \renewcommand\chaptername{Chapter}
}
\selectlanguage{english}

% Headings in Times bold (manual §8: Arial or Times everywhere). Without this,
% abnTeX2 requests TeX Gyre Heros (qhv) and falls back to Computer Modern.
\renewcommand{\ABNTEXchapterfont}{\rmfamily\bfseries}
\renewcommand{\ABNTEXsectionfont}{\rmfamily\bfseries}
\renewcommand{\ABNTEXsubsectionfont}{\rmfamily\bfseries}
\renewcommand{\ABNTEXsubsubsectionfont}{\rmfamily\bfseries}

\hyphenation{fe-de-ral}
\hyphenation{scientiae}

% ---------------------------------------------------------------------------------------
% THE SUPPLEMENTARY VOLUME, named in one place. [round7, 2026-07-29]
% Appendices B (errata) and D (the label-history benchmark) live in a separate document,
% src/main_extra.tex, per the author's decision on the advisor's ruling of 2026-07-28. Eight
% sentences in this document point the reader at that volume, and they must all name it the
% same way: a reader who meets "the supplementary volume", "the supplementary material" and
% "the extra volume" in three chapters has to work out whether those are one thing.
% \extravolume is that one name. Change it here and every pointer follows.
% NOT a \ref: the target is in another document and has no label here. The pointers say what
% the volume is called, not what number anything in it carries -- except the two appendix
% letters B and D, which are stable by construction (0_extra.tex sets the chapter counter so
% they keep the letters they had here) and are guarded by src_utils/check_extra_xrefs.py.
\newcommand{\extravolume}{the supplementary volume of this dissertation}
\newcommand{\Extravolume}{The supplementary volume of this dissertation}
% ---------------------------------------------------------------------------------------

% Chapter preface: the time-capsule device (NORTH_STAR §3). One italic paragraph
% before §N.1 stating venue, status, contribution note, and what later chapters revise.
\newenvironment{chapterpreface}%
{\begin{quotation}\itshape\noindent}%
{\end{quotation}\vspace{0.5\baselineskip}}

% ---------------------------------------------------------------
% Cover / folha de rosto metadata
% ---------------------------------------------------------------
% TITLE (author decision 2026-07-24): the working title is now SELECTED and live below.
% The Resumo/Abstract catalog headers echo it verbatim (kept in sync). This remains a
% "for now" choice pending the final call with the advisor; the alternates stay commented
% here so the author can swap in one line if the advisor prefers another.
% Echo locations kept in sync: \titulo below (folha de rosto), Resumo catalog header,
% Abstract catalog header; hyperref pdftitle inherits \titulo.
%
% SELECTED (author 2026-07-24):
%   From Representations to a Single Joint Model: Multi-Task Learning for
%   Point-of-Interest Category and Region Prediction
%
% ALTERNATES (kept for the advisor conversation; all pass the two-factor rule -- name
% both the representation and the sharing/architecture factor, or neither;
% storyline/AVAL_NECESSARIA_3_ptBR.md D2):
%   1. Multi-Task Learning for Point-of-Interest Classification and Prediction Tasks: The Role of the Check-in-Level Representation
%   2. Representation Drives, Sharing Delivers: A Check-in-Level Multi-Task
%      Study of Next Category and Region Prediction
%   3. The Visit, Not the Place: Check-in-Level Representations and a Joint
%      Model for Next Category and Region Prediction
%   4. One Visit, One Vector: From a Published Null to a Joint Model for
%      Next Category and Region Prediction
\titulo{\normalsize{\textbf{Multi-Task Learning for Point-of-Interest Classification and Prediction Tasks: The Role of the Check-in-Level Representation}}}
\autor{\normalsize{\textbf{Vitor Hugo De Oliveira Silva}}}
\local{\normalsize{\textbf{Florestal - Minas Gerais}}}
\data{\normalsize{\textbf{2026}}}
\orientador{Fabrício Aguiar Silva}
\instituicao{Universidade Federal de Viçosa}
% \campus set 2026-07-27 (author: "o campus e o Florestal", src_utils/PENDENCIAS.md:37).
% Consumed by \imprimircampus in abntex2-UFV.sty:49, inside \imprimircapa only; that macro
% is not called by either build today, so this renders nothing until a cover page is added.
\campus{Campus Florestal}
\curso{Pós-graduação em Ciência da Computação}
\membrobancaA{[Banca member 1 --- pending advisor conversation]}
\membrobancaB{[Banca member 2 --- pending advisor conversation]}
\databanca{[defense date --- pending]}

\preambulo{Dissertation submitted to the Computer Science Graduate Program of the
Universidade Federal de Viçosa in partial fulfillment of the requirements for the
degree of \textit{Magister Scientiae}.}

\makeatletter
\hypersetup{
    pdftitle={\@title},
    pdfauthor={\@author},
    pdfsubject={\imprimirpreambulo},
    pdfcreator={LaTeX with abnTeX2},
    colorlinks=true,
    linkcolor=black,
    citecolor=blue,
    filecolor=magenta,
    urlcolor=blue,
    bookmarksdepth=4,
    % [round6, E-5, visual/presentation persona 18] hyperfootnotes DISABLED.
    % THE DEFECT: every build logged ten "pdfTeX warning (dest): name{Hfootnote.N} has been
    % referenced but does not exist" lines, and the persona measured where the links actually
    % LAND rather than only whether the destinations resolve: all ten footnote marks jumped the
    % reader to PHYSICAL PAGE 1 (the folha de rosto in the defense build, the List of Figures in
    % the deposit build). Ten dead links in a PDF handed to a committee, and the counter that
    % would have shown them is one no build claim in this project has ever carried.
    % THE CAUSE: abntex2 loads hyperref itself, and this document's footnotes are set inside
    % groups the footnote-anchor mechanism does not survive, so hyperref emits a reference to an
    % anchor it never writes. The link then falls back to the start of the document.
    % THE FIX: turn the footnote marks back into plain text. A footnote mark sits on the same
    % page as its own note in this document, so the link was never carrying the reader anywhere
    % useful even when it worked. Nothing else in hyperref changes: internal cross-references,
    % citations and URLs stay live.
    % NOTE: the option is NOT set here. Setting hyperfootnotes=false in \hypersetup after
    % abntex2 has already loaded hyperref leaves all ten warnings in place (measured). It is
    % passed at load time with \PassOptionsToPackage before \documentclass, at the top of this
    % file. This comment stays here because this is where someone will look for it.
}
\makeatother

